\section{Related Work}
\label{sec:related}

\paragraph{Lexicon-based detection.}
Pre-transformer detectors used lexicons and hand-written surface features
\cite{waseem2016hateful}: Gitari et al.\ \cite{gitari2015lexicon} score a comment
from a curated profanity lexicon and apply rules over the matched phrases---inspectable,
but hand-curated, with recall bounded by the list. Wiegand et al.\
\cite{wiegand2018inducing} automate that induction step, the closest antecedent to how
our entries acquire severity and confidence. Schmidt and Wiegand's survey
\cite{schmidt2017survey} explains the move to learned models: the \hate{}/\off{}
boundary is combinatorial, not lexical. Davidson et al.\ \cite{davidson2017automated}
found the two classes share most of their surface vocabulary---the tension that makes
our channel over-fire on \off{}---and surveys \cite{fortuna2018survey} list the
recurring difficulties: imbalance, subjectivity, overlap, obfuscation, domain shift.
We take these as the object of analysis rather than as obstacles to a classifier.

\paragraph{Vietnamese resources and models.}
ViHSD \cite{luu2021vihsd} is the reference three-way benchmark, ViCTSD
\cite{nguyen2021victsd} provides related data, and ViCM \cite{nguyen2026vicm} adds
code-mixed comments. PhoBERT \cite{nguyen2020phobert}
adapts RoBERTa to Vietnamese, ViSoBERT \cite{nguyen2023visobert} is pre-trained on
social media, and ViHateT5 \cite{nguyen2024vihatet5} casts the task as generation.
ViHOS \cite{hoang2023vihos} annotates the \emph{spans} that make a comment hateful,
the closest Vietnamese work to our notion of evidence; our spans come from a lexicon
matcher rather than human annotation, but the motivation is the same. SMTCE
\cite{nguyen2022smtce} standardises Vietnamese social-media evaluation, which is why we
report macro-F1 rather than accuracy alone.

\paragraph{Multilingual transformers.}
mBERT \cite{devlin2019bert} and XLM-R \cite{conneau2020xlmr} established multilingual
pre-training; DeBERTa \cite{he2021deberta} introduced disentangled attention and
DeBERTa-v3 \cite{he2023debertav3} adds ELECTRA-style pre-training, whose
\texttt{microsoft/\allowbreak mdeberta-v3-base} checkpoint \cite{liu2024mdeberta} we
fine-tune.

\paragraph{Explainability and hybrid systems.}
HateXplain \cite{mathew2021hatexplain} made human rationales and explanation quality
benchmark quantities. LIME \cite{ribeiro2016why} and SHAP
\cite{lundberg2017shap} are model-agnostic but do not guarantee that the highlighted
tokens are the ones the decision used. Our rule channel is part of the decision
function, so its matched spans are \emph{faithful by construction} for the lexical
component, while the neural component remains opaque. We differ in two ways: we fix
the accuracy question with a paired test and analyse provenance instead, and we
report the rule channel's own error profile and the direction of every rule-induced
change. To our knowledge, this per-decision provenance audit has not been reported
for Vietnamese hate speech.
